% =====================================================================
%  config/commands.tex
%  Perintah kustom yang dipakai bersama oleh banyak bab.
%  Sebelum menambah perintah baru, periksa dulu apakah fungsinya sudah ada
%  di sini (Aturan AI #2: jangan menduplikasi konfigurasi).
% =====================================================================

% --- Himpunan bilangan ---------------------------------------------------
\newcommand{\R}{\mathbb{R}}
\newcommand{\N}{\mathbb{N}}
\newcommand{\Z}{\mathbb{Z}}

% --- Penanda teks --------------------------------------------------------
\newcommand{\code}[1]{\texttt{#1}}
\newcommand{\important}[1]{\textbf{Penting:} #1}

% --- Keterangan sumber gambar -------------------------------------------
% Dipakai di bawah gambar yang diambil dari bahan kuliah, agar atribusinya
% seragam dan tidak perlu ditulis ulang di setiap bab.
\newcommand{\sumber}[1]{\par\vspace{2pt}{\footnotesize\itshape Sumber: #1\par}}

% --- Path berkas kode program -------------------------------------------
% Berkas kode disimpan di code/<bahasa>/ (Arsitektur_Buku_Ajar.md Bagian 19)
% dan ditarik dengan \lstinputlisting. Berbeda dari \input dan
% \includegraphics, \lstinputlisting tidak diperbaiki path-nya oleh paket
% subfiles, sehingga path relatifnya pecah pada salah satu mode kompilasi:
% kompilasi penuh berjalan dari buku_ajar/ (butuh "code/..."), sedangkan
% kompilasi mandiri bab berjalan dari chapters/NN/ (butuh "../../code/...").
% Perintah ini memilih bentuk yang benar pada kedua mode.
\newcommand{\codefile}[1]{%
    \ifSubfilesClassLoaded{../../code/#1}{code/#1}%
}
